\chapter{Contrôle stochastique}\label{ch:qm:stochastic-control}

Après une baisse de 20~\% des actions, la pondération en actions d'un fonds de pension est passée de 60~\% à 52~\%. Son
règlement interne prévoit~: rééquilibrer à 60~\% en fin de trimestre, ce qui revient à acheter des actions après une baisse. La moitié
du comité d'investissement veut attendre que le marché se stabilise. La question relève du contrôle
stochastique~: choisir, à chaque instant et avec ce qui est connu alors, la quantité à détenir, de manière à maximiser
l'espérance de l'utilité de sa richesse finale. Pour un investisseur à \omterm{def:qm:stochastic-control:utility}{aversion relative au risque} constante dans un
marché à prix lognormaux, Merton l'a résolu en 1969~: détenir une fraction constante, quoi qu'il soit arrivé,
et rééquilibrer vers elle. Avec les hypothèses de ce fonds, la fraction vaut 51~\%, alors que détenir 60~\% coûte
3.6~points de base de rendement en \omterm{def:qm:stochastic-control:ce}{équivalent certain} par an, et le débat du comité porte sur un chiffre
bien plus petit que l'incertitude affectant son propre rendement espéré. Ce chapitre développe la \omterm{def:qm:stochastic-control:dp}{programmation
dynamique} en temps discret, l'\omterm{def:qm:stochastic-control:hjb}{équation de Hamilton--Jacobi--Bellman} en temps continu, le
théorème de vérification qui transforme un candidat en preuve, le problème de Merton et une esquisse des
\omterm{def:qm:stochastic-control:viscosity}{solutions de viscosité} qui traitent les \omterm{def:qm:stochastic-control:problem}{fonctions valeur} présentant des points anguleux.

\section{Programmation dynamique}

\begin{definition}[Problème de contrôle stochastique, contrôle admissible, fonction valeur]\label{def:qm:stochastic-control:problem}
Un \emph{problème de contrôle stochastique}\index{problème de contrôle stochastique} choisit un processus $u_t$ qui
influence la dynamique d'un état $X$, ici $dX_t = \mu(t, X_t, u_t)\,dt + \sigma(t, X_t, u_t)\,dW_t$ ou son
analogue en temps discret, pour maximiser $J(t, x; u) = \E_{t,x}[\int_t^Tf(s, X_s, u_s)\,ds + g(X_T)]$. Un
\emph{contrôle admissible}\index{contrôle admissible} est un \omterm{def:qm:probability-at-speed:filtration}{processus adapté} à valeurs dans un ensemble $\mathcal U$ pour
lequel l'équation d'état admet une solution et $J$ est bien définie. La \emph{fonction valeur}\index{fonction valeur}
est $V(t, x) = \sup_uJ(t, x; u)$ sur les contrôles admissibles.
\end{definition}

\begin{theorem}[Principe de la programmation dynamique]\label{thm:qm:stochastic-control:dpp}
Pour tout $t < s \le T$, $V(t, x) = \sup_u\E_{t,x}\bigl[\int_t^sf(r, X_r, u_r)\,dr + V(s, X_s)\bigr]$.
\end{theorem}

\begin{proof}[Preuve partielle]
En temps discret, avec un nombre fini d'états et de contrôles~: toute stratégie issue de $t$ engrange ses gains jusqu'à
$s$, puis au plus $V(s, X_s)$, ce qui donne $\le$~; suivre une continuation optimale (ou $\varepsilon$-optimale)
issue de chaque état en $s$, après une première partie quelconque, donne $\ge$. En temps continu, l'argument
requiert une sélection mesurable de continuations $\varepsilon$-optimales (Fleming et Soner, 2006).
\end{proof}

\begin{definition}[Programmation dynamique, contrôle en boucle fermée]\label{def:qm:stochastic-control:dp}
La \emph{programmation dynamique}\index{programmation dynamique} calcule la \omterm{def:qm:stochastic-control:problem}{fonction valeur} à rebours dans le temps
à l'aide du principe à un pas, $V_t(x) = \max_u\{f(t, x, u) + \E[V_{t+1}(X_{t+1}) \mid X_t = x, u]\}$, à partir de $V_T =
g$. Un \emph{contrôle en boucle fermée}\index{contrôle en boucle fermée} est un contrôle de la forme $u_t = u^*(t, X_t)$, fonction de
l'état courant~; la programmation dynamique en produit un.
\end{definition}

Le principe ramène un choix portant sur des stratégies entières à une suite de choix à un pas, au prix de
calculer $V$ sur tout l'espace des états~: la malédiction de la dimension. Le solveur du projet courant
(le premier listing du tutoriel) le fait sur une grille, l'espérance étant calculée par
quadrature sur l'état suivant et interpolation linéaire entre les points de la grille.

\section{L'équation de Hamilton--Jacobi--Bellman}

\begin{definition}[Équation de Hamilton--Jacobi--Bellman]\label{def:qm:stochastic-control:hjb}
Pour la diffusion contrôlée de la \cref{def:qm:stochastic-control:problem}, l'\emph{équation de
Hamilton--Jacobi--Bellman}\index{équation de Hamilton--Jacobi--Bellman} est
\[
\partial_tV(t, x) + \sup_{u\in\mathcal U}\Bigl\{\mu(t, x, u)\partial_xV + \tfrac12\sigma^2(t, x, u)\partial_{xx}V + f(t, x,
u)\Bigr\} = 0, \qquad V(T, x) = g(x).
\]
\end{definition}

C'est le principe de la \omterm{def:qm:stochastic-control:dp}{programmation dynamique} sur un intervalle $[t, t + h]$ avec $h \to 0$~: développer $V(t + h,
X_{t+h})$ par la formule d'Itô, prendre les espérances, diviser par $h$. Pour chaque contrôle fixé, c'est l'équation rétrograde
du chapitre~4, avec le générateur $\mathcal L^u$ de l'état contrôlé~; le supremum la rend non linéaire.
L'heuristique suppose $V$ régulière, ce que le théorème de vérification remplace par une
hypothèse portant sur un candidat.

\section{Vérification}

\begin{theorem}[Vérification]\label{thm:qm:stochastic-control:verification}
Soit $w \in C^{1,2}$ à croissance polynomiale solution de l'\omterm{def:qm:stochastic-control:hjb}{équation de Hamilton--Jacobi--Bellman}, et supposons que le supremum soit atteint en
$u^*(t, x)$ tel que l'équation d'état, avec le \omterm{def:qm:stochastic-control:dp}{contrôle en boucle fermée} $u^*$, admette une solution. Alors $w = V$ et $u^*$ est
optimal, à condition que les intégrales stochastiques ci-dessous soient de vraies \omterm{def:qm:probability-at-speed:martingale}{martingales}.
\end{theorem}

\begin{proof}
Pour tout $u$ admissible, la formule d'Itô donne $w(T, X_T) = w(t, x) + \int_t^T(\partial_s + \mathcal L^{u_s})w\,ds + \int_t^T\sigma
\partial_xw\,dW_s$. L'\omterm{def:qm:stochastic-control:hjb}{équation de Hamilton--Jacobi--Bellman} affirme que $(\partial_s + \mathcal L^u)w + f \le 0$, donc, en prenant les espérances, $\E[g(X_T) + \int_t^Tf\,ds] \le
w(t, x)$~: $J \le w$. Avec $u = u^*$, l'inégalité est une égalité, donc $J(u^*) = w \ge V \ge J(u^*)$.
\end{proof}

\section{Problème de Merton}

\begin{definition}[Fonction d'utilité, aversion relative au risque, utilité CRRA]\label{def:qm:stochastic-control:utility}
Une \emph{fonction d'utilité}\index{fonction d'utilité} $U$ est croissante et concave~; l'investisseur maximise
$\E[U(W_T)]$. Son \emph{aversion relative au risque}\index{aversion relative au risque} est $-wU^{\prime\prime}(w)/U'(w)$.
Une \emph{utilité CRRA}\index{utilité CRRA} a une aversion relative au risque constante $\gamma > 0$~: $U(w) =
w^{1-\gamma}/(1 - \gamma)$, et $\ln w$ pour $\gamma = 1$.
\end{definition}

\begin{definition}[Équivalent certain]\label{def:qm:stochastic-control:ce}
L'\emph{équivalent certain}\index{équivalent certain} d'une richesse aléatoire $W$ est le montant certain $c$
avec $U(c) = \E[U(W)]$~; sur un horizon $T$, le rendement en équivalent certain vaut $\ln(c/W_0)/T$.
\end{definition}

Un fonds détient une fraction $\pi_t$ de sa richesse dans un actif risqué avec $dS/S = \mu\,dt + \sigma\,dW$ et le reste dans le
\omterm{def:qm:girsanov-and-changes-of-numeraire:numeraire}{compte monétaire} au taux $r$, de sorte que $dW_t = W_t(r + \pi_t(\mu - r))\,dt + W_t\pi_t\sigma\,dW_t$.

\begin{theorem}[Merton]\label{thm:qm:stochastic-control:merton}
Pour une \omterm{def:qm:stochastic-control:utility}{utilité CRRA} sur la richesse terminale, la fraction optimale est constante,
\[
\pi^* = \frac{\mu - r}{\gamma\sigma^2},
\]
et la \omterm{def:qm:stochastic-control:problem}{fonction valeur} vaut $V(t, w) = \frac{w^{1-\gamma}}{1-\gamma}e^{(1-\gamma)\rho(T-t)}$, où $\rho = r + \frac{(\mu -
r)^2}{2\gamma\sigma^2}$ est le rendement en \omterm{def:qm:stochastic-control:ce}{équivalent certain} de la stratégie optimale.
\end{theorem}

\begin{proof}
On essaie $V(t, w) = \frac{w^{1-\gamma}}{1-\gamma}h(t)$ dans l'\omterm{def:qm:stochastic-control:hjb}{équation de Hamilton--Jacobi--Bellman}
\[
\partial_tV + \sup_\pi\bigl\{w(r + \pi(\mu - r))\partial_wV + \tfrac12w^2\pi^2\sigma^2\partial_{ww}V\bigr\} = 0.
\] Comme $w\partial_wV = (1 - \gamma)V$ et $w^2\partial_{ww}V = -\gamma(1 - \gamma)V$, l'expression entre accolades vaut
$(1 - \gamma)V(r + \pi(\mu - r) - \tfrac12\gamma\pi^2\sigma^2)$, maximale en $\pi^*$ avec pour valeur $(1 - \gamma)V\rho$. D'où $h' +
(1 - \gamma)\rho h = 0$, $h(T) = 1$. Le candidat est régulier et l'équation de la richesse à $\pi^*$ constant est un
\omterm{def:qm:stochastic-differential-equations:gbm}{mouvement brownien géométrique}, donc le \cref{thm:qm:stochastic-control:verification} s'applique.
\end{proof}

\begin{definition}[fraction de Merton]\label{def:qm:stochastic-control:fraction}
La \emph{fraction de Merton}\index{fraction de Merton} est $\pi^* = (\mu - r)/(\gamma\sigma^2)$.
\end{definition}

La fraction optimale ne dépend ni de la richesse ni du temps~: après une baisse, le fonds doit racheter des actions pour revenir à cette fraction.
Détenir une fraction $\pi$ constante rapporte le rendement d'\omterm{def:qm:stochastic-control:ce}{équivalent certain} $r + \pi(\mu - r) - \tfrac12\gamma\pi^2\sigma^2$
(la \cref{fig:qm:stochastic-control:ce}), donc une fraction erronée coûte $\tfrac12\gamma\sigma^2(\pi - \pi^*)^2$ par an. Avec $\mu =
7\%$, $r = 2\%$, $\sigma = 18\%$ et $\gamma = 3$~: $\pi^* = 51.4\%$, un rendement d'\omterm{def:qm:stochastic-control:ce}{équivalent certain} de 3.29~\%, et 60~\%
coûte 3.6~points de base par an. Avec $\gamma = 1$, utilité logarithmique, $\pi^* = (\mu - r)/\sigma^2$ est la fraction de Kelly
du livre~2, chapitre~29, ici 154~\%~: l'investisseur optimal pour la croissance emprunte.

\begin{omfigure}
\begin{tikzpicture}
\begin{axis}[width=11cm,height=4.8cm,
  xlabel={fraction en actions $\pi$},ylabel={rendement d'équivalent certain (\% par an)},
  tick label style={font=\small},label style={font=\small},
  xmin=0,xmax=1.6,ymin=1.5,ymax=3.5,grid=major,grid style={gray!25},
  table/col sep=comma]
\addplot[omDef,very thick] table[x=pi,y=ce]{figdata/methods/09-stochastic-control/ce.csv};
\addplot[black,only marks,mark=*,mark size=1.8pt] coordinates {(0.5144,3.2860) (0.60,3.2504)};
\addplot[gray,dashed] coordinates {(1.5432,1.5) (1.5432,3.5)};
\node[font=\footnotesize,anchor=east] at (axis cs:1.52,2.9) {Kelly ($\gamma = 1$)};
\end{axis}
\end{tikzpicture}

\omcaption{Rendement d'\omterm{def:qm:stochastic-control:ce}{équivalent certain} d'une fraction constante en actions pour un investisseur à \omterm{def:qm:stochastic-control:utility}{utilité CRRA} de $\gamma = 3$
($\mu = 7\%$, $r = 2\%$, $\sigma = 18\%$)~: une parabole culminant à la \omterm{def:qm:stochastic-control:fraction}{fraction de Merton} (point de gauche, 51.4~\%), assez plate pour que 60~\% (point de droite) ne coûte que
3.6~pb par an. La fraction de Kelly, optimale pour $\gamma = 1$, vaut 154~\%. Données~: forme fermée, le
tutoriel du chapitre.}\label{fig:qm:stochastic-control:ce}
\end{omfigure}

La platitude du sommet a un revers. La fraction est proportionnelle au rendement excédentaire espéré, qui est
connu bien moins précisément~: estimé sur dix ans de rendements à 18~\% de volatilité, $\mu$ a une erreur type de $18\%/\sqrt{10}
= 5.7\%$, et chaque point de $\mu$ déplace $\pi^*$ de dix points (la \cref{fig:qm:stochastic-control:sensitivity}).
La règle de Merton est exacte, mais son entrée est incertaine~; le chapitre~14 réduit cette entrée.

\begin{omfigure}
\begin{tikzpicture}
\begin{groupplot}[group style={group size=2 by 1,horizontal sep=1.6cm},
  width=7.6cm,height=4.8cm,tick label style={font=\small},label style={font=\small},
  grid=major,grid style={gray!25},table/col sep=comma]
\nextgroupplot[xlabel={rendement excédentaire espéré $\mu - r$ (\%)},ylabel={fraction de Merton},xmin=0,xmax=10,ymin=0,ymax=1.6,
  legend style={font=\footnotesize,at={(0.03,0.97)},anchor=north west,fill=white,draw=none}]
\addplot[omThm,very thick] table[x=excess,y=g2]{figdata/methods/09-stochastic-control/sensitivity.csv};
\addplot[omDef,very thick] table[x=excess,y=g3]{figdata/methods/09-stochastic-control/sensitivity.csv};
\addplot[omProp,very thick] table[x=excess,y=g5]{figdata/methods/09-stochastic-control/sensitivity.csv};
\legend{$\gamma = 2$,$\gamma = 3$,$\gamma = 5$}
\nextgroupplot[xlabel={années},ylabel={poids en actions},xmin=0,xmax=10,ymin=0.4,ymax=0.9]
\addplot[name path=hi,gray!40,forget plot] table[x=year,y=p90]{figdata/methods/09-stochastic-control/drift.csv};
\addplot[name path=lo,gray!40,forget plot] table[x=year,y=p10]{figdata/methods/09-stochastic-control/drift.csv};
\addplot[gray!20,forget plot] fill between[of=hi and lo];
\addplot[omDef,very thick] table[x=year,y=p50]{figdata/methods/09-stochastic-control/drift.csv};
\addplot[black,dashed,domain=0:10] {0.5144};
\end{groupplot}
\end{tikzpicture}

\omcaption{À gauche~: la \omterm{def:qm:stochastic-control:fraction}{fraction de Merton} en fonction du rendement excédentaire espéré, pour trois aversions au risque ($\sigma = 18\%$)~:
dix points de fraction par point de rendement à $\gamma = 3$. À droite~: le poids en actions d'un fonds qui part de 60~\% et
ne se rééquilibre jamais, médiane et bande des 10\up{e}--90\up{e} centiles sur 20\,000 trajectoires simulées, comparé à la \omterm{def:qm:stochastic-control:fraction}{fraction de Merton}
(en pointillés). Données~: le tutoriel du chapitre, à graine fixée.}\label{fig:qm:stochastic-control:sensitivity}
\end{omfigure}

\section{Esquisse des solutions de viscosité}

Les \omterm{def:qm:stochastic-control:problem}{fonctions valeur} ne sont souvent pas régulières~: un contrôle qui bascule d'un extrême à l'autre à une frontière crée
un coude dans $V$, et l'équation HJB ne peut pas être vérifiée en ce point au sens classique. La bonne notion remplace les dérivées par
des fonctions test régulières qui touchent $V$.

\begin{definition}[solution de viscosité]\label{def:qm:stochastic-control:viscosity}
Une fonction continue $V$ est une \emph{solution de viscosité}\index{solution de viscosité} de $F(x, V, DV, D^2V) = 0$ ($F$ non croissante
en son dernier argument) si, en tout point où une fonction régulière $\phi$ touche $V$ par au-dessus ($V - \phi$ admet un maximum local),
$F(x, V, D\phi, D^2\phi) \le 0$, et en tout point où elle touche par en dessous, $F \ge 0$.
\end{definition}

Pour le problème de sortie de $(-1, 1)$ aussi vite que possible à vitesse unitaire, $V(x) = 1 - |x|$ et $|V'| = 1$ partout sauf
en zéro, où aucune dérivée n'existe~; c'est la \omterm{def:qm:stochastic-control:viscosity}{solution de viscosité}, et elle est la limite, pour $\varepsilon \to 0$, des
solutions régulières de $-\varepsilon u^{\prime\prime} + |u'| = 1$ (la \cref{fig:qm:stochastic-control:viscosity}), d'où vient le
nom. Crandall et Lions (1983) ont démontré l'existence et l'unicité pour les équations du premier ordre~; la théorie
garantit que la \omterm{def:qm:stochastic-control:dp}{programmation dynamique} sur grille converge vers la bonne \omterm{def:qm:stochastic-control:problem}{fonction valeur}, même lorsqu'elle présente des coudes, comme dans les
problèmes d'arrêt du chapitre~10.

\begin{omfigure}
\begin{tikzpicture}
\begin{axis}[width=10cm,height=4.6cm,
  xlabel={$x$},ylabel={$u(x)$},
  tick label style={font=\small},label style={font=\small},
  xmin=-1,xmax=1,ymin=0,ymax=1.1,grid=major,grid style={gray!25},
  legend columns=3,legend style={font=\footnotesize,at={(0.5,-0.3)},anchor=north,draw=none,fill=none,/tikz/every even column/.append style={column sep=8pt}},
  table/col sep=comma]
\addplot[omProp,very thick] table[x=x,y=e02]{figdata/methods/09-stochastic-control/viscosity.csv};
\addplot[omThm,very thick] table[x=x,y=e005]{figdata/methods/09-stochastic-control/viscosity.csv};
\addplot[omDef,very thick,dashed] table[x=x,y=e0]{figdata/methods/09-stochastic-control/viscosity.csv};
\legend{$\varepsilon = 0.2$,$\varepsilon = 0.05$,$1 - |x|$}
\end{axis}
\end{tikzpicture}

\omcaption{Solutions de $-\varepsilon u^{\prime\prime} + |u'| = 1$ sur $(-1, 1)$ avec $u(\pm1) = 0$, à savoir $1 - |x| + \varepsilon(e^{-1/\varepsilon} -
e^{-|x|/\varepsilon})$~: régulières pour tout $\varepsilon > 0$, qui convergent vers la \omterm{def:qm:stochastic-control:problem}{fonction valeur} à coude $1 - |x|$, \omterm{def:qm:stochastic-control:viscosity}{solution
de viscosité} de $|u'| = 1$.}\label{fig:qm:stochastic-control:viscosity}
\end{omfigure}

\section{Tutoriel~: Merton par induction rétrograde}\label{tut:qm:stochastic-control:merton}

\begin{tutorial}
\textbf{Objectif.} Résoudre une version à temps discret du problème du fonds par \omterm{def:qm:stochastic-control:dp}{programmation dynamique}, la confronter
à la \omterm{def:qm:stochastic-control:fraction}{fraction de Merton}, et mesurer ce que rapporte le rééquilibrage. \textbf{État final~:}
les \cref{fig:qm:stochastic-control:ce,fig:qm:stochastic-control:sensitivity} et les chiffres du problème du week-end.
\begin{tutsteps}
\item \textbf{Le solveur}~: induction rétrograde sur une grille, avec une récompense d'arrêt optionnelle pour le chapitre~10.
\omcode{code/firm/dpsolve/firm_dpsolve.py}{30}{49}{Induction rétrograde avec quadrature et interpolation.}\label{lst:qm:stochastic-control:dp}
\item \textbf{Le problème}~: dix dates annuelles de rééquilibrage, logarithme de la richesse sur une grille de 241 points, fractions de 0 à 1.5
par pas de 0.005, rendements annuels lognormaux par quadrature de Gauss--Hermite à 20 points.
\omcode{code/methods/09-stochastic-control/python/qm_merton.py}{66}{84}{Problème de Merton traité comme programmation dynamique sur grille.}\label{lst:qm:stochastic-control:merton}
\item \textbf{Exécuter} \texttt{problem()}~: la fraction optimale vaut 0.52 à tout niveau de richesse et à toute date, le point de grille
le plus proche de l'optimum à une période, 0.517 (le rééquilibrage continu donne 0.514)~; puis \texttt{fig\_merton.py}.
\end{tutsteps}
\textbf{Pour aller plus loin.} Ajouter une consommation à fraction constante de la richesse et vérifier la règle de consommation de Merton~;
remplacer l'\omterm{def:qm:stochastic-control:utility}{utilité CRRA} par une utilité avec plancher (un passif que le fonds doit couvrir) et observer la fraction commencer à dépendre de
la richesse.
\end{tutorial}

\section{Construction~: le solveur de programmation dynamique}\label{bld:qm:stochastic-control:dpsolve}

\begin{build}
\bfield{Rôle} Résoudre les petits problèmes de contrôle de la maison (allocation, inventaire, calendriers d'exécution) et, avec une
récompense d'arrêt, ses problèmes d'arrêt (chapitre~10), par induction rétrograde sur une grille.
\bfield{Interface} \texttt{backward\_induction(grid, controls, n\_steps, reward, transition, terminal, stop=None)}
renvoyant les valeurs, la politique et la région d'arrêt~; \texttt{gauss\_hermite\_normal(n)}.
\bfield{Règles} Grille d'états unidimensionnelle, ensemble fini de contrôles, espérances calculées à l'aide de nœuds et de poids fournis par l'utilisateur, interpolation
linéaire, extrapolation constante~; vectorisée sur les états et les contrôles.
\bfield{Tests d'acceptation} \texttt{code/firm/dpsolve/tests/}~: la quadrature de Gauss--Hermite intègre exactement les polynômes~; le problème
de Merton renvoie une fraction indépendante de la richesse, égale à l'optimum à une période~; un problème d'arrêt de
solution connue (la sortie de type américain du chapitre~10) est reproduite.
\bfield{Pour aller plus loin} États à deux dimensions avec interpolation bilinéaire~; itération sur les politiques pour les horizons infinis.
\end{build}

\begin{omsources}
\item R.~C.~Merton, «~Lifetime portfolio selection under uncertainty~: the continuous-time case~», \emph{Review of
Economics and Statistics} 51, 1969.
\item P.~A.~Samuelson, «~Lifetime portfolio selection by dynamic stochastic programming~», \emph{Review of Economics and
Statistics} 51, 1969.
\item R.~Bellman, \emph{Dynamic Programming}, Princeton University Press, 1957.
\item M.~G.~Crandall et P.-L.~Lions, «~Viscosity solutions of Hamilton--Jacobi equations~», \emph{Transactions of the
AMS} 277, 1983.
\item W.~H.~Fleming et H.~M.~Soner, \emph{Controlled Markov Processes and Viscosity Solutions}, Springer, 2\up{e} éd., 2006.
\end{omsources}

\section{Exercices}

\begin{exercise}[$\star$]\label{exo:qm:stochastic-control:1}
Un actif présente un rendement excédentaire espéré de 4~\% et une volatilité de 20~\%. Déterminer la fraction qu'un investisseur
d'\omterm{def:qm:stochastic-control:utility}{aversion relative au risque} $\gamma = 4$ détient~?
\end{exercise}

\begin{exercise}[$\star$]\label{exo:qm:stochastic-control:2}
Avec les hypothèses du fonds, calculer le rendement en \omterm{def:qm:stochastic-control:ce}{équivalent certain} d'une détention de 100~\% en actions~?
\end{exercise}

\begin{exercise}[$\star$]\label{exo:qm:stochastic-control:3}
Calculer l'\omterm{def:qm:stochastic-control:utility}{aversion relative au risque} de $U(w) = \ln w$ et de $U(w) = w^{1-\gamma}/(1 - \gamma)$.
\end{exercise}

\begin{exercise}[$\star\star$]\label{exo:qm:stochastic-control:4}
Montrer que détenir une fraction constante $\pi$ au lieu de $\pi^*$ coûte $\tfrac12\gamma\sigma^2(\pi - \pi^*)^2$ de rendement en \omterm{def:qm:stochastic-control:ce}{équivalent certain} par an,
et évaluer ce coût pour 60~\% contre 51.4~\%.
\end{exercise}

\begin{exercise}[$\star\star$]\label{exo:qm:stochastic-control:5}
Avec les hypothèses du fonds, calculer la fraction détenue par un investisseur à utilité logarithmique. Expliquer pourquoi un fonds ne le ferait pas.
\end{exercise}

\begin{exercise}[$\star\star$]\label{exo:qm:stochastic-control:6}
Vérifier que $V = \frac{w^{1-\gamma}}{1-\gamma}e^{(1-\gamma)\rho(T-t)}$ résout l'équation HJB, avec $\rho = r + (\mu - r)^2/(2\gamma\sigma^2)$, et
calculer $\rho$ pour le fonds.
\end{exercise}

\begin{exercise}[$\star\star\star$]\label{exo:qm:stochastic-control:7}
\emph{Programmation.} Avec \texttt{backward\_induction}, résoudre le problème du fonds avec $\gamma = 5$ et un rééquilibrage annuel. Vérifier si
la fraction reste indépendante de la richesse, et la comparer à $(\mu - r)/(\gamma\sigma^2)$.
\end{exercise}

\begin{exercise}[$\star\star\star$]\label{exo:qm:stochastic-control:8}
\emph{Repérer l'erreur.} «~Nous avons estimé la prime de risque des actions à 5~\% à partir de dix ans de données, donc notre pondération optimale en actions est
51~\%, et nous conserverons exactement ce niveau.~»
\end{exercise}

\section{Problème~: le comité de rééquilibrage}

\begin{problem}[{Problème du week-end --- ce que vaut la décision du comité}]\label{pb:qm:stochastic-control:1}
Un fonds à \omterm{def:qm:stochastic-control:utility}{utilité CRRA}, $\gamma = 3$, anticipe un rendement des actions de $\mu = 7\%$ avec une volatilité $\sigma = 18\%$~; les liquidités rapportent $r =
2\%$. Sa politique d'investissement fixe 60~\% d'actions, avec rééquilibrage. Après une baisse, sa pondération tombe à 52~\%.

\medskip\noindent\textbf{Partie I --- Merton.}
\begin{enumerate}
\item Quelle est la \omterm{def:qm:stochastic-control:fraction}{fraction de Merton}~?
\item Quel est le rendement en \omterm{def:qm:stochastic-control:ce}{équivalent certain} de la stratégie optimale~?
\item Combien coûte par an une détention de 60~\%~?
\item Et une détention de 52~\%~?
\item Quelle transaction Merton recommande-t-il après la baisse, et que recommande le règlement de marché~?
\end{enumerate}

\medskip\noindent\textbf{Partie II --- Sensibilité.}
\begin{enumerate}[resume]
\item Quelles sont les fractions pour $\gamma = 2$ et $\gamma = 5$~?
\item Et pour $\mu = 6\%$ et $8\%$~?
\item Que donne $\gamma = 1$, et que signifie ce cas~?
\item Que donnent la \omterm{def:qm:stochastic-control:dp}{programmation dynamique} sur grille avec rééquilibrage annuel et l'optimum à une période~?
\item Pourquoi la fraction ne dépend-elle pas de la richesse~?
\end{enumerate}

\medskip\noindent\textbf{Partie III --- Rééquilibrage.}
\begin{enumerate}[resume]
\item Sur dix années simulées, quels rendements en \omterm{def:qm:stochastic-control:ce}{équivalent certain} obtiennent une fraction constante de 60~\%, la \omterm{def:qm:stochastic-control:fraction}{fraction de Merton} et une stratégie buy-and-hold partie de 60~\%~?
\item Où se situe la pondération buy-and-hold au terme de dix ans (10\up{e} centile, médiane, 90\up{e})~?
\item Que coûte le fait de ne jamais rééquilibrer, comparé à un rééquilibrage à 60~\%~?
\item Avec quelle erreur type $\mu$ est-il connu à partir de dix ans de rendements, et quelle fourchette cela donne-t-il pour $\pi^*$~?
\item Lequel compte le plus pour le fonds~: le choix du comité entre 52~\% et 60~\%, ou l'incertitude sur $\mu$~?
\end{enumerate}

\medskip\noindent\textbf{Partie IV --- Jugement.}
\begin{enumerate}[resume]
\item Que changeraient les coûts de transaction~?
\item Que changeraient les passifs (une promesse de retraite)~?
\item Que doit dire le directeur des investissements au comité~?
\item Énoncer le \emph{résultat nommé}~: la \omterm{def:qm:stochastic-control:fraction}{fraction de Merton} et le coût de la fraction de 60~\% prévue par le règlement de marché.
\item En une phrase~: que dit la solution de Merton sur le rééquilibrage après une baisse~?
\end{enumerate}
\end{problem}

\section{Questions d'entretien}

\begin{interviewq}[$\star$ \iqroles{chercheur, trader}]\label{iq:qm:stochastic-control:1}
Énoncez la fraction optimale de Merton et expliquez chacun de ses trois paramètres.
\end{interviewq}

\begin{interviewq}[$\star\star$ \iqroles{chercheur}]\label{iq:qm:stochastic-control:2}
Dérivez l'équation HJB à partir du principe de \omterm{def:qm:stochastic-control:dp}{programmation dynamique}.
\end{interviewq}

\begin{interviewq}[$\star\star$ \iqroles{chercheur}]\label{iq:qm:stochastic-control:3}
Pourquoi un investisseur CRRA, dans le modèle de Merton, achète-t-il des actions après leur baisse~?
\end{interviewq}

\begin{interviewq}[$\star\star$ \iqroles{trader, chercheur}]\label{iq:qm:stochastic-control:4}
Comment le critère de Kelly se rattache-t-il au problème de Merton~?
\end{interviewq}

\begin{interviewq}[$\star\star$ \iqroles{développeur}]\label{iq:qm:stochastic-control:5}
Votre \omterm{def:qm:stochastic-control:dp}{programmation dynamique} sur grille renvoie une politique qui saute entre des points voisins de la grille. Que vérifiez-vous~?
\end{interviewq}

\begin{interviewq}[$\star\star\star$ \iqroles{chercheur}]\label{iq:qm:stochastic-control:6}
Quand la \omterm{def:qm:stochastic-control:problem}{fonction valeur} d'un problème de contrôle n'est-elle pas régulière, et comment donne-t-on un sens à l'équation HJB dans ce cas~?
\end{interviewq}
